\subsection{R-19: simultaneous row and column blocks improve the range again}

\paragraph{Statement R-19.}
For $m\geq10$ put
\[
 a_m=\min\{\lfloor m/3\rfloor,m-9\},\qquad
 H_m=4a_m(m-9).
\]
Then
\begin{equation}\label{R19:dimension}
 \dim\mathcal P_{m,4}\geq2m^2+2+H_m.
\end{equation}
If $n>m\geq10$ and
\begin{equation}\label{R19:padding}
 2n^2\leq2m^2+H_m+2,
\end{equation}
the four-row multiplicity gaps of R-14/G-10 and the explicit
rectangular gaps of G-11 both hold. In particular, they hold
eventually for $n=\lfloor\rho m\rfloor$ and every fixed
$1<\rho<\sqrt{5/3}$. Earlier lower bounds remain available
when they are larger at a particular finite size.

\begin{proof}
Use the same cyclic pencil with
$\ell_i=x-\zeta^iy$, $v_i=z-\zeta^iw$.
Choose centers
\[
 \mathcal A=\{1,4,\ldots,1+3(a_m-1)\}.
\]
The triples of rows $a-1,a,a+1$ are disjoint and any two
centers have cyclic distance at least three.
For each center add
\begin{equation}\label{R19:threechords}
 p_a=\varepsilon\ell_{a+1}\text{ at }(a,a+2),\qquad
 q_a=\varepsilon v_{a+2}\text{ at }(a+1,a),\qquad
 r_a=\varepsilon\ell_a\text{ at }(a-1,a+1).
\end{equation}
The two rows $a,a+1$ still have support only in columns
$a,a+1,a+2$. The two columns $a,a+1$ have support only
in rows $a-1,a,a+1$. To verify the latter assertion for
all simultaneous motifs, an additional forward chord
from another center $b$ could enter these columns only
if $b=a-2$ or $a-1$, and an additional backward or third
chord only if $b=a-1$ or $a+1$. These cases are excluded
by the spacing. The analogous row-support assertion is
immediate from the disjoint row triples.

Thus the two row identities of R-18 remain exact and
give the same families $U_{a,k}(\varepsilon)$ and
$V_{a,k}(\varepsilon)$, with the remaining complementary
permanents now evaluated at the three-chord pencil.
There are also two exact column identities. For a column
$c\notin\{a,a+1\}$ let $S_{a,c}(\varepsilon)$ be the
permanent obtained by deleting rows $a-1,a,a+1$ and
columns $a,a+1,c$. The complementary permanents factor as
\[
 \begin{aligned}
 C_{a-1,c}&=(\ell_a\ell_{a+1}+v_aq_a)S_{a,c},\\
 C_{a,c}&=(v_{a-1}\ell_{a+1}+q_ar_a)S_{a,c},\\
 C_{a+1,c}&=(v_{a-1}v_a+\ell_ar_a)S_{a,c}.
 \end{aligned}
\]
Indeed the two remaining rows of the triple must match
the two supported columns; all other matching choices
occur in the same complementary matrix.
Consequently
\begin{align*}
 v_{a-1}C_{a-1,c}-\ell_a C_{a,c}-q_a C_{a+1,c}
    &=-2q_ar_a\ell_a S_{a,c},\\
 r_a C_{a-1,c}+v_a C_{a,c}-\ell_{a+1}C_{a+1,c}
    &= 2q_ar_av_a S_{a,c}.
\end{align*}
After division by the nonzero scalar $\varepsilon^2$,
these show that the differential image contains
\[
 C_{a,k}(\varepsilon)=\ell_a^2v_{a+2}S_{a,a-k}(\varepsilon),
 \qquad
 D_{a,k}(\varepsilon)=\ell_av_av_{a+2}S_{a,a-k+1}(\varepsilon).
\]
We use $5\leq k\leq m-5$ below.
At zero the unique cycle matching gives
\[
 \begin{aligned}
 C_{a,k}&=\ell_a^2v_{a+2}
       \prod_{i=a-k}^{a-2}v_i
       \prod_{i\notin[a-k,a+1]}\ell_i,\\
 D_{a,k}&=\ell_av_av_{a+2}
       \prod_{i=a-k+1}^{a-2}v_i
       \prod_{i\notin[a-k+1,a+1]}\ell_i.
 \end{aligned}
\]
All four families $U,V,C,D$ have bidegree $(m-k,k)$,
and translating their center by $a$ applies $\tau^a$.

We prove fourfold normal independence in each character
\begin{equation}\label{R19:interior}
 5\leq k,t\leq m-5.
\end{equation}
As in R-18 write $q=\zeta$, $K=q^k$, $X=q^t$, and
divide character-$t$ coordinates by the nonzero
coefficients $A_{t-s}$ of
$A(x,y)=\prod_{i=k}^{m-1}(x-q^iy)$.
Set
\[
 J(z,w)=\prod_{i=2}^{k-4}(z-q^iw),\qquad
 v_b=z-q^bw.
\]
The two base columns and the two row-family columns,
after this normalization, are $J$ times
$v_{k-3}v_{k-2}$ times the four independent cubics
of \eqref{R18:cubics}. Thus these four quintics span
\begin{equation}\label{R19:rowquintics}
 v_{k-3}v_{k-2}\operatorname{Sym}^3\mathbb C^2.
\end{equation}

We compute the two column-family sequences explicitly.
Put $M=m-k$ and
$\widetilde A=A(x,q^{2-k}y)=\prod_{i=2}^{M+1}(x-q^iy)$.
Their $(x,y)$ factors at center zero are
\[
 A_C=(x-y)^2\prod_{i=2}^{M-1}(x-q^iy),\qquad
 A_D=(x-y)\prod_{i=2}^{M}(x-q^iy).
\]
Define
\[
 \begin{aligned}
 \delta_0&=\frac{K(q-1)}{q(K-1)},&
 \delta_1&=\frac{K-q}{q(K-1)},\\
 \eta_0&=\frac{K^2(q-1)^2}{q(K-1)(qK-1)},&
 \eta_1&=\frac{K(K-q)(q-1)}{q(K-1)(qK-1)},&
 \eta_2&=\frac{K-1}{qK-1}.
 \end{aligned}
\]
The following are polynomial identities:
\[
 \begin{aligned}
 A_D&=\delta_0\widetilde A(x,y)
           +\delta_1\widetilde A(x,q^{-1}y),\\
 A_C&=\eta_0\widetilde A(x,y)
           +\eta_1\widetilde A(x,q^{-1}y)
           +\eta_2\widetilde A(x,q^{-2}y).
 \end{aligned}
\]
For the first, cancel the common factors with indices
$2,\ldots,M$ and compare the remaining linear forms.
For the second cancel those with indices $2,\ldots,M-1$;
the remaining identity is
\[
 \eta_0(x-K^{-1}y)(x-qK^{-1}y)
 +\eta_1(x-qy)(x-K^{-1}y)
 +\eta_2(x-y)(x-qy)=(x-y)^2,
\]
which follows by its three coefficients. This verifies
the coefficient ratios without any genericity assumption.

Since $\widetilde A_r=q^{(2-k)r}A_r$,
the two column vectors, after the common coordinate
division and multiplication by the nonzero character
scalar $X^{k-2}$, are $J$ times the quintics
\begin{align*}
 \mathcal C_X={}&
 \eta_0v_{-2}v_{-1}v_0v_1v_k
 +\eta_1X^{-1}v_{-1}v_0v_1v_{k-3}v_{k+1}
 +\eta_2X^{-2}v_0v_1v_{k-3}v_{k-2}v_{k+2},\\
 \mathcal D_X={}&
 \delta_0v_{-1}v_0v_1v_{k-2}v_k
 +\delta_1X^{-1}v_0v_1v_{k-3}v_{k-1}v_{k+1}.
\end{align*}
For example the $(z,w)$ factor of $C_{0,k}$ has roots
$2,m-k,\ldots,m-2$; shifting $w$ by $q^{k-2}$ changes
them to $k,-2,-1,0,\ldots,k-4$.
Its next two shifts give the other two terms above.
The factor of $D_{0,k}$ has roots
$0,2,m-k+1,\ldots,m-2$; its corresponding two shifts
give the terms of $\mathcal D_X$. This checks every
factor and character scalar in these expressions.

To test independence from \eqref{R19:rowquintics},
evaluate the two new quintics at
$z/w=K/q^3$ and $K/q^2$, the distinct roots of
$v_{k-3}v_{k-2}$. Substitution gives
\begin{equation}\label{R19:evaldet}
\begin{aligned}
 &\begin{vmatrix}
 \mathcal C_X(K/q^3,1)&\mathcal D_X(K/q^3,1)\\
 \mathcal C_X(K/q^2,1)&\mathcal D_X(K/q^2,1)
 \end{vmatrix}\\
 &\qquad=
 \frac{K^6(q-K)(q^2-K)^2(q^3-K)^2(q^4-K)
             (q-1)^6(q+1)(q^2+q+1)}
      {q^{27}(K-1)(qK-1)}\ne0.
\end{aligned}
\end{equation}
Here too the cancellation of $X$ is exact.
Every displayed factor is nonzero: $m\geq10$ and
$5\leq k\leq m-5$, so $q$ has order neither two nor
three and $K$ differs from $1,q,q^2,q^3,q^4,q^{-1}$.
The four quintics in \eqref{R19:rowquintics} and the
two new ones consequently form a basis of all binary
quintics.

Under \eqref{R19:interior}, the allowed indices $s$
contain six consecutive values $b,\ldots,b+5$, where
$0\leq b\leq k-5$. Multiplication by $J$ and restriction
to these six coordinates has, in the quintic monomial
basis, determinant
\[
 \det(J_{b+u-v})_{0\leq u,v\leq5}
   =\pm s_{(6^b)}(q^2,q^3,\ldots,q^{k-4})\ne0.
\]
There are $N=k-5$ variables. The Weyl numerator
exponents for $(6^b)$ are distinct in
$[0,N+5]=[0,k]$, below the order $m$, so the same
Vandermonde proof as in R-16 applies. When $N=0$,
the determinant is one. Thus all six normalized
character columns are independent; equivalently,
the four classes of $U,V,C,D$ are independent modulo
the base character space.

A relation among their translates at centers in
$\mathcal A$ now gives four separate equations
$\sum_{a\in\mathcal A}c_aq^{at}=0$ in each of the
available characters. Since $a_m\leq m-9$, the
$a_m$ consecutive characters
$t=5,\ldots,a_m+4$ lie in the allowed range.
Each system is an invertible Vandermonde system in
the distinct $q^a$. Hence the $4a_m$ classes are
independent in each bidegree. Different bidegrees
are independent, giving $4a_m(m-9)=H_m$ classes
modulo $W_0$.

All four families have polynomial deformations lying
in the actual differential image for nonzero
$\varepsilon$. The augmented-minor argument of
R-16, with the $2m^2+2$ base columns, proves
\eqref{R19:dimension}.
R-14 and G-10 give the general four-row consequence.
Furthermore $H_m\leq(4m/3)(m-9)$, so
\eqref{R19:padding} implies
$n^2\leq5m^2/3-6m+1<16m^2/9$.
Thus $n<4m/3$ and G-11 also applies.
Finally $H_m=4m^2/3+O(m)$, giving the asserted
asymptotic ratio $\sqrt{5/3}$.
\end{proof}

